% ============================================================================
%  Part VII: Distributed Systems
% ============================================================================
\section{Distributed Systems}
\label{sec:distributed}

This part covers the questions that arise when many machines generate TesseraIDs concurrently and consumers need to make sense of the resulting stream.

\subsection{Clock skew under the Standard variant}

TesseraID generators on different machines are subject to whatever clock skew exists between those machines. NTP \citep{rfc5905} can synchronize machine clocks to roughly 1 ms in a controlled data center and 10--100 ms across the open internet; chronyd or PTP \citep{ieee1588} can do better in tightly-controlled environments. TesseraID's behavior under skew depends on the variant.

\textbf{Standard variant.} Two TesseraIDs generated on machines whose clocks differ by 50 ms will have timestamps that differ by 50 ms. Sorting them puts the earlier-clocked machine first, which may or may not reflect the actual order of the underlying events. Within a single machine, ordering is exact (the counter resolves intra-millisecond ties); across machines, ordering reflects clock skew.

For most use cases this is fine. Database indexes don't care about cross-node ordering accuracy; they care about each node's own writes being clustered. Logs aggregated from multiple machines are usually re-ordered by an external time-sorting layer that explicitly accounts for clock skew. If you need finer cross-node ordering than NTP delivers, use the HLC variant.

\subsection{HLC: causal ordering across nodes}

The HLC variant solves the cross-node ordering problem when nodes exchange messages. The mechanism is simple: every message between nodes carries the sender's HLC value, and every receiver \emph{observes} that HLC before generating its next ID. The observer's HLC then incorporates the maximum of (its own HLC, the peer's HLC, the local wall clock); the result is that any ID generated after observing a peer's message has an HLC value at least as large as the peer's.

The practical pattern:

\begin{lstlisting}[language=Pseudocode,caption={Plumbing HLC observation through a message-passing layer},label={lst:hlc-plumbing}]
# Sender side
msg.hlc <- generator.current_hlc()  # snapshot of (l, c)
network.send(msg)

# Receiver side
network.receive(msg)
generator.observe(msg.hlc)          # advance our HLC if msg's is larger
local_id <- generator.new_hlc()     # next ID respects causal order
\end{lstlisting}

The 60-second drift bound rejects observations that would push logical time more than $\Delta_{\max}$ ahead of the local physical clock. This is the safety valve that prevents a malicious or buggy peer from advancing the local HLC arbitrarily far. Sixty seconds is large enough to absorb cross-region NTP skew and short network outages, but small enough that a clock that drifts further surfaces as a hard error rather than silent corruption.

\subsection{When HLC and Standard TesseraIDs mix}

A consumer that sorts a stream of mixed Standard and HLC TesseraIDs gets correct lexicographic byte order. What it does \emph{not} get is meaningful semantic ordering --- a Standard TesseraID's timestamp is wall-clock-relative, while an HLC TesseraID's timestamp may be ahead of wall-clock by up to 60 seconds depending on observed peer clocks. Mixing them in a single ordered stream is permitted by the format but introduces an ordering-by-format-type artifact that consumers should be aware of.

Best practice: choose one variant per logical stream. A given table, log, or event channel should be all-Standard or all-HLC, not a mix. If a system genuinely needs both, keep them in separate streams.

\subsection{Coordination and uniqueness}

TesseraID is uncoordinated: there is no central authority assigning machine IDs or ranges to generators. Uniqueness comes from the random tail. The probability that two independent generators produce the same TesseraID is bounded by the birthday math: at $10^9$ IDs/sec generated by a million-node cluster, the expected time to first collision is $1.35 \times 10^{22}$ years.

Compared to coordinated formats like Snowflake \citep{snowflake} (which require every generator to have a pre-assigned unique ID and break catastrophically if the assignment is wrong), TesseraID's uncoordinated design is more robust to operational errors. New nodes can come up at any time without registration; failed nodes can be replaced without re-issuing identifiers.

The cost is byte count: 41 bytes vs 8 for Snowflake. Snowflake's compactness comes from compressing what TesseraID puts in the random tail (cryptographic uniqueness) into a coordination service (machine-ID assignment correctness). For systems that already have such a service, Snowflake is unbeatable on size; for systems that don't, TesseraID is unbeatable on robustness.

\subsection{Clock skew detection and alarming}

A practical operational concern: if a node's clock drifts more than the HLC drift bound, its peers start rejecting its messages. This is the desired behavior --- silent corruption is worse than loud failure --- but it means operators need to monitor for the failure mode.

The recommended practice is to log every \code{HLC\_drift\_exceeded} error with the offending logical timestamp, the local physical timestamp, and the peer node identity. Alerting fires when the rate of such errors crosses a threshold. The log makes it easy to diagnose which clock drifted in which direction.

\subsection{Snapshots and consistent cuts}

For systems that need to take consistent global snapshots (e.g., for backup, compliance, or analytics), the HLC variant provides a natural ``snapshot at logical time $T$'' operation: include every event whose HLC is $\leq T$. Because HLC preserves causality, this snapshot is causally consistent: if event A causes event B, the snapshot either includes both or excludes both, never just B.

This is the same guarantee Chandy--Lamport \citep{chandylamport1985} provides via marker messages, but achieved without the explicit marker-passing protocol. The HLC TesseraID itself is the marker.

\subsection{Cross-region replication}

For multi-region deployments, the HLC variant simplifies replication ordering. A row written in region A and replicated to region B carries its HLC in the row metadata. When B applies the replication, it observes A's HLC; subsequent writes in B inherit causal ordering from A. The result is a globally-sortable stream of writes that respects causality across regions, even when wall clocks across regions disagree.

The Standard variant does not give you this. Replicating Standard TesseraIDs across regions produces a stream that sorts by wall-clock time, which may or may not reflect the order of the underlying events.

